\documentclass{article}

\usepackage[italian]{babel} %testi autogenerati italiano
\usepackage{minted}         %per codice vhdl bello
\usepackage{tikz}           %per disegno fsm
\usetikzlibrary{automata, positioning, arrows}
\usepackage{circuitikz}     %per disegno componente
\usepackage{graphicx}       %per importare immagini
\usepackage{geometry}       %per gestire margini e spostamenti
\geometry {
    top=25mm,
    bmargin=25mm,
}
\usepackage{array}          %per colonne di width fissata
\usepackage{subcaption}     %tabelle divise

\begin{document}

\setlength\parindent{0pt} %noindent automatico

\begin{titlepage}
    \centering
    \hrule

    \vspace{0,5cm}
    {
        \normalsize Politecnico di Milano\\
	                Dipartimento di Elettronica, Informazione e Bioingegneria \par
    }
    
    \vspace{4,5cm}
    {\Huge \textbf{Progetto di Reti Logiche\\
    2020/21}\\}

    \vspace{0,5cm}
    \large {Prof. Palermo Gianluca}
        
    \vspace{2,5cm}
    {
        \large
        \begin{tabular}{c c}
            Shalby Hazem Hesham Yousef & (Codice Persona: 10596243, Matricola: 910871)\\
            Perego Niccolò & (Codice Persona: 10628782, Matricola: 895468)\\
        \end{tabular}
    
    }
    
    \vspace{5,2cm}

        \normalsize{1 Aprile 2021 \par}
        \vspace{0,3cm}

        \centering\hspace{0,2cm}\includegraphics[scale=0.6]{logo.png}
        \vspace{0,5cm}
        \hrule

\end{titlepage}

\pagebreak

\section{Requisiti di progetto} %1

\subsection{Descrizione del problema} %1.1
Si vuole realizzare un componente in grado di svolgere una versione semplificata del processo di equalizzazione dell’istogramma di un’immagine, ossia di ricalibrare il contrasto di quest’ultima, \\effettuando una ridistribuzione dei valori di intensità pixel per pixel. \\\\
Le immagini di cui è richiesta la manipolazione sono definite in scala di grigi a 256 livelli e hanno una dimensione massima di 128x128 pixel.

\vspace{0,5cm} %un po' di spazio

\subsection{Interfaccia del componente} %1.2
Il componente realizzato ha un’interfaccia così definita in liguaggio VHDL:

\begin{minted}{vhdl}
    ENTITY project_reti_logiche IS PORT
	(
		i_clk     : IN  std_logic;
		i_rst     : IN  std_logic; 
		i_start   : IN  std_logic; 
		i_data    : IN  std_logic_vector (7 DOWNTO 0);
		o_address : OUT std_logic_vector (15 DOWNTO 0);
		o_done    : OUT std_logic; 
		o_en      : OUT std_logic;
		o_we      : OUT std_logic;
		o_data    : OUT std_logic_vector (7 DOWNTO 0) 
	);
    END project_reti_logiche;
\end{minted}

\vspace{0,5cm} %un po' di spazio

In particolare:
\begin{itemize}
    \item \texttt{i\_clk} è il segnale di CLOCK in ingresso generato dal TestBench;
    \item \texttt{i\_rst} è il segnale di RESET che inizializza la macchina, predisponendola alla ricezione del segnale di START. Può essere anche asincrono;
    \item \texttt{i\_start} è il segnale di START generato dal Test Bench;
    \item \texttt{i\_data} è il segnale (vettore) che arriva dalla memoria in seguito a una richiesta di lettura;
    \item \texttt{o\_address} è il segnale (vettore) di uscita che manda l’indirizzo alla memoria;
    \item \texttt{o\_done} è il segnale di uscita che comunica la fine dell’elaborazione e il dato di uscita scritto in memoria;
    \item \texttt{o\_en} è il segnale di ENABLE da mandare alla memoria per poter comunicare (sia in lettura che in scrittura);
    \item \texttt{o\_we} è il segnale di WRITE ENABLE da mandare alla memoria per comunicare quale operazione si vuole svolgere su di essa. Può assumere valori 0 e 1, rispettivamente per lettura e scrittura;
    \item \texttt{o\_data} è il segnale (vettore) di uscita dal componente verso la memoria.
\end{itemize}

\pagebreak

\begin{figure}[ht]

    \centering
    \begin{circuitikz}
        \ctikzset{multipoles/thickness=3.5}
        \ctikzset{multipoles/dipchip/width=4}
        %sinistra
        \draw   (0,0) node[dipchip,
                num pins=10, hide numbers, no topmark,
                external pins width=0](C){\large{\texttt{COMPONENT}}};
        \node   [right] at (C.bpin 1) {\texttt{i\_start}};
        \draw   (C.bpin 1) -- ++(-1,0) coordinate(extpin);
        \node   [right] at (C.bpin 2) {\texttt{i\_data}};
        \draw   (C.bpin 2) -- ++(-1,0) coordinate(extpin);
        \node   [right] at (C.bpin 4) {\texttt{i\_rst}};
        \draw   (C.bpin 4) -- ++(-1,0) coordinate(extpin);
        \draw   (C.bpin 5) ++(0,0.2) -- ++(0.2,-0.2)
                node[right]{\texttt{i\_clk}} -- ++(-0.2,-0.2);
        \draw   (C.bpin 5) -- ++(-1,0) coordinate(extpin);
        %destra
        \node   [left] at (C.bpin 6) {\texttt{o\_data}};
        \draw   (C.bpin 6) -- ++(1,0) coordinate(extpin);
        \node   [left] at (C.bpin 7) {\texttt{o\_done}};
        \draw   (C.bpin 7) -- ++(1,0) coordinate(extpin);
        \node   [left] at (C.bpin 8) {\texttt{o\_en}};
        \draw   (C.bpin 8) -- ++(1,0) coordinate(extpin);
        \node   [left] at (C.bpin 9) {\texttt{o\_we}};
        \draw   (C.bpin 9) -- ++(1,0) coordinate(extpin);
        \node   [left] at (C.bpin 10) {\texttt{o\_address}};
        \draw   (C.bpin 10) -- ++(1,0) coordinate(extpin);
        \end{circuitikz}

    \setlength{\belowcaptionskip}{0.1cm} %un po' di spazio
    \caption{Schema del componente realizzato.}
    \label{fig:component}
\end{figure}

\subsection{Descrizione della memoria e dell'interazione con il componente} %1.3
Il modulo implementato dovrà dialogare in lettura e scrittura con una RAM, indirizzata al byte.\\\\
In particolare, l’algoritmo di equalizzazione sarà applicato a immagini pre-salvate in memoria, la cui grandezza effettiva (in pixel) sarà specificata dal prodotto del contenuto tra le celle a indirizzo \texttt{0} e \texttt{1} della RAM, contenenti rispettivamente il numero di colonne \texttt{n\_col} e di righe \texttt{n\_row} dell’immagine, entrambi di dimensione 8 bit.\\
Nei byte successivi, dall’indirizzo \texttt{2} all’indirizzo \texttt{n\_col $\cdot$ n\_row $+$ 1}, sarà contenuta, pixel per pixel, sequenzialmente e in modo contiguo, l’immagine di cui è richiesta la trasformazione.\\
Infine, l’immagine ottenuta dal processo di equalizzazione svolto dal componente verrà salvata in memoria dall’indirizzo \texttt{n\_col $\cdot$ n\_row $+$ 2} all’indirizzo \texttt{2 $\cdot$ n\_col $\cdot$ n\_row $+$ 1}.

\vspace{0,5cm} %un po' di spazio

\begin{table}[h]
    \centering
    \def\arraystretch{1.3} %un po' di padding

    \caption{Schema generale del contenuto della memoria dopo l'elaborazione.}

    \begin{tabular}{ |m{4cm} | m{5cm} }
        \cline{1-1}
        \texttt{n\_col} & \texttt{addr. 0} \\
        \cline{1-1}
        \texttt{n\_row} & \texttt{addr. 1} \\
        \cline{1-1}
        \texttt{Primo pixel immagine da elaborare} & \texttt{addr. 2} \\
        \cline{1-1}
        \vdots & \vdots \\
        \cline{1-1}
        \texttt{Ultimo pixel immagine da elaborare} & \texttt{addr. n\_col $\cdot$ n\_row $+$ 1} \\
        \cline{1-1}
        \texttt{Primo pixel immagine elaborata} & \texttt{addr. n\_col $\cdot$ n\_row $+$ 2} \\
        \cline{1-1}
        \vdots & \vdots \\
        \cline{1-1}
        \texttt{Ultimo pixel immagine elaborata} & \texttt{addr. 2 $\cdot$ n\_col $\cdot$ n\_row $+$ 1} \\
        \cline{1-1}    
        \end{tabular}
    
    \label{tab:schemamem}
    
\end{table}

\vspace{0.3cm}

\subsection{Esempio di funzionamento} %1.4
Si riporta in seguito un'esempio di elaborazione compiuta su un'immagine di test.

\pagebreak

\begin{figure}[ht]
    \hspace{-10pt}
    \begin{minipage}[c]{.35\linewidth}
    \centering 
    \begin{tabular}{c|c}
        addr. & data \\
        \hline \hline
        0  & 4   \\
        1  & 4   \\
        2  & 52  \\
        3  & 62  \\
        4  & 85  \\
        5  & 87  \\
        6  & 55  \\
        7  & 59  \\
        8  & 71  \\
        9  & 78  \\
        10 & 122 \\
        11 & 68  \\
        12 & 59  \\
        13 & 68  \\
        14 & 70  \\
        15 & 113 \\
        16 & 64  \\
       \end{tabular}
    \begin{tabular}{c|c}
        addr. & data \\
        \hline \hline
        17 & 69  \\
        18 & 0   \\
        19 & 12  \\
        20 & 255 \\
        21 & 72  \\
        22 & 40  \\
        23 & 28  \\
        24 & 64  \\
        25 & 244 \\
        26 & 132 \\
        27 & 76  \\
        28 & 28  \\
        29 & 48  \\
        30 & 140 \\
        31 & 104 \\
        32 & 64  \\
        33 & 68  \\
       \end{tabular}
    \subcaption{Contenuto della memoria}\label{esempio1}
    \end{minipage}%
    \begin{minipage}[c]{.3\linewidth}
    \centering
        \includegraphics[scale=0.25]{immSorg.jpg}
    \subcaption{Immagine sorgente}\label{esempio2}
    \end{minipage}\hspace{20pt}
    \begin{minipage}[c]{.3\linewidth}
        \centering
            \includegraphics[scale=0.25]{immElab.jpg}
        \subcaption{Immagine equalizzata}\label{esempio3}
        \end{minipage}
    \end{figure}

\vspace{0.3cm}

È evidente come l'immagine \small (c) \normalsize presenti un contrasto maggiore rispetto alla \small(b)\normalsize. Questo è dato dall'ampliamento del range di valori assunti dai pixel dell'immagine, come evidenziato dalla tabella di memoria:
\begin{itemize}
    \item Immagine sorgente, addr. 2 - 17: i pixel assumono valori da 52 a 122;
    \item Immagine equalizzata, addr. 18 - 33: i pixel assumono valori da 0 a 255.
\end{itemize}

Non avendo ancora affrontato appieno il processo svolto dal componente, si ritiene opportuno non approfondire in questa sezione 
l'insieme di passaggi che permettono la trasformazione dell'immagine riportata nell'esempio, che verranno ripresi in seguito. 
\\\\
Si rende disponibile un TestBench che replica l'elaborazione dell'immagine \small (b) \normalsize a questo link:\\

\end{document}
